\chapter{Relativistic Plasma Radiation and Spectral Source Terms}
\label{ch:plasma-radiation}

\section{Kinetic description of a relativistic electron population}

For an isotropic relativistic electron population, momentum is conveniently normalized as
\begin{equation}
 u \equiv \frac{p}{m_ec}, \qquad \gamma=\sqrt{1+u^2}, \qquad E=(\gamma-1)m_ec^2.
\end{equation}
The dimensionless temperature is
\begin{equation}
 \theta \equiv \frac{k_B T_e}{m_ec^2}.
\end{equation}
When thermal equilibrium is a defensible approximation, the Maxwell--J\"uttner distribution can be written
\begin{equation}
 f_{\mathrm{MJ}}(\mathbf{p})
 =\frac{n_e}{4\pi m_e^3c^3\theta K_2(1/\theta)}
 \exp\!\left(-\frac{\gamma}{\theta}\right),
\label{eq:mj}
\end{equation}
with normalization $\int f_{\mathrm{MJ}}\,\dd^3p=n_e$. The corresponding probability density in $\gamma$ is proportional to $\gamma^2\beta\exp(-\gamma/\theta)$, where $\beta=v/c$.

Equation~\eqref{eq:mj} is a reference state, not a guarantee. In a driven plasma the stationary distribution solves a kinetic equation schematically written
\begin{equation}
\frac{\partial f}{\partial t}
+\dot{\mathbf{x}}\cdot\nabla f
+\dot{\mathbf{p}}\cdot\nabla_p f
=C_{ee}[f]+C_{ei}[f]+Q-S-L,
\label{eq:kinetic-general}
\end{equation}
where $C_{ee}$ and $C_{ei}$ are collision operators, $Q$ represents injection or acceleration, $S$ represents continuous sinks such as radiation reaction, and $L$ represents orbit or boundary loss. The proposed magnetic-control strategy acts mainly through $L$ and, indirectly, through changes to $\dot{\mathbf{x}}$ and $\dot{\mathbf{p}}$.

\section{Classical radiation background}

An accelerated charge radiates. In the nonrelativistic limit, the Larmor power is
\begin{equation}
 P_{\mathrm{L}}=\frac{e^2a^2}{6\pi\epsilon_0c^3}.
\label{eq:larmor}
\end{equation}
The relativistic Liénard generalization is
\begin{equation}
P=\frac{e^2\gamma^6}{6\pi\epsilon_0c^3}
\left[a^2-(\boldsymbol{\beta}\times\mathbf{a})^2\right],
\label{eq:lienard}
\end{equation}
which emphasizes the strong energy dependence of radiation from relativistic electrons \cite{jackson1999}. Bremsstrahlung is the quantum-mechanical realization of radiation during Coulomb scattering. The photon emission probability is of order the fine-structure constant relative to elastic scattering, but the enormous number of encounters in a dense plasma makes the integrated power important \cite{jauch1976,akhiezer1965,bethe1954,koch1959}.

\section{Electron--ion bremsstrahlung}

For an ion species $i$ with number density $n_i$ and charge $Z_i e$, the differential spectral emissivity is
\begin{equation}
 j_{ei}(E_\gamma)
 =n_i\int_{E_\gamma}^{\infty}
 v(E)f_E(E)
 E_\gamma\frac{\dd\sigma_{ei}(E,E_\gamma;Z_i)}{\dd E_\gamma}
 \,\dd E,
\label{eq:ei-emissivity}
\end{equation}
where $f_E(E)\dd E$ is the electron number density in the interval $[E,E+\dd E]$. The total power density is
\begin{equation}
 P_{ei}=\sum_i\int_0^\infty j_{ei}(E_\gamma)\,\dd E_\gamma.
\label{eq:pei}
\end{equation}
The Bethe--Heitler treatment, with Coulomb and screening corrections appropriate to the regime, supplies the differential cross section \cite{bethe1954,davies1954,koch1959,haug2004}. The ion contribution scales roughly with $\sum_i n_iZ_i^2$, commonly represented using
\begin{equation}
 Z_{\mathrm{eff}}=\frac{\sum_i n_iZ_i^2}{n_e}.
\end{equation}
This weighting makes impurities important: a trace high-$Z$ species can increase radiation substantially even if it contributes little to the total particle density.

In a nonrelativistic, fully ionized plasma, the familiar frequency-integrated scaling is
\begin{equation}
P_{ei}^{\mathrm{NR}}\propto n_en_iZ_i^2T_e^{1/2}\overline{g}_{\mathrm{ff}},
\label{eq:ei-nr}
\end{equation}
where $\overline{g}_{\mathrm{ff}}$ is the thermally averaged Gaunt factor. Relativistic corrections increase the high-temperature emission and change the spectrum. The simple $T_e^{1/2}$ scaling should therefore not be extrapolated into the hundreds-of-keV regime without correction \cite{nozawa1998,svennson1982}.

\section{Electron--electron bremsstrahlung}

Electron--electron bremsstrahlung is suppressed at low energy because equal charge-to-mass particles have no changing electric dipole moment in their center-of-momentum frame. At relativistic energy, higher multipoles and kinematic effects make the channel significant. A general spectral expression is
\begin{equation}
 j_{ee}(E_\gamma)
 =\frac{1}{2}\int\!\dd^3p_1\int\!\dd^3p_2\,
 f(\mathbf{p}_1)f(\mathbf{p}_2)
 v_{\mathrm{rel}}E_\gamma
 \frac{\dd\sigma_{ee}}{\dd E_\gamma},
\label{eq:ee-emissivity}
\end{equation}
with the factor $1/2$ preventing double counting. The calculation is more expensive than the electron--ion integral because it depends on two electron momenta and their relative angle. In relativistic fusion studies, however, omission can bias the total loss and spectral shape \cite{haug1975,svennson1982,munirov2023}.

A frequently used fit for fully ionized relativistic plasmas separates electron--ion and electron--electron terms as functions of $\theta$ and $Z_{\mathrm{eff}}$. Such fits are useful for power-balance benchmarks, but the present thesis requires differential spectra because biological dose conversion is energy dependent. Integrated fits should therefore be used to verify, not replace, the spectral calculation.

\section{Synchrotron and cyclotron radiation}

A magnetized electron radiates because the Lorentz force bends its trajectory. For an ultrarelativistic electron with pitch angle $\alpha$, the synchrotron power in vacuum is
\begin{equation}
P_{\mathrm{syn}}
=\frac{e^4B^2\gamma^2\beta^2\sin^2\alpha}
{6\pi\epsilon_0m_e^2c^3}.
\label{eq:synch-power}
\end{equation}
The characteristic angular frequency scales as
\begin{equation}
\omega_c=\frac{3}{2}\gamma^2\omega_{ce}\sin\alpha,
\qquad \omega_{ce}=\frac{eB}{m_e}.
\label{eq:synch-freq}
\end{equation}
Synchrotron power therefore competes directly with any strategy that changes pitch angle, energy, or magnetic-field strength. A perturbation that reduces bremsstrahlung by moving electrons to large perpendicular momentum could unintentionally increase synchrotron emission. Conversely, selective loss of large-$p_\perp$ electrons may reduce synchrotron at the cost of anisotropy or current changes. Bremsstrahlung and synchrotron must be optimized jointly \cite{blumenthal1970,embreus2016}.

\section{Inverse Compton and other photon channels}

Energetic electrons can upscatter ambient photons by inverse Compton scattering. In the Thomson limit, the average scattered photon energy grows approximately as $\gamma^2$, while the power per electron is
\begin{equation}
P_{\mathrm{IC}}=\frac{4}{3}\sigma_Tc\gamma^2\beta^2U_{\mathrm{rad}},
\end{equation}
where $U_{\mathrm{rad}}$ is the photon energy density. In dense laboratory fusion plasmas, inverse Compton power is often below bremsstrahlung and synchrotron, but it may matter in laser-driven, externally illuminated, or radiation-trapping configurations. Pair production becomes possible for photons above \SI{1.022}{\MeV}; photonuclear interactions and activation require still more source-specific treatment. These channels are retained in the hazard inventory even when not dominant in the baseline calculation.

\section{Energy-cutoff and redistribution models}

A hard truncation
\begin{equation}
 f_{\mathrm{tr}}(E)=A f_{\mathrm{MJ}}(E)H(E_c-E)
\label{eq:hardcut}
\end{equation}
removes particles above cutoff $E_c$. If density is held fixed by choosing $A$, the distribution is renormalized but particle energy and entropy are altered. If the physical mechanism is escape, the removed density must be replenished by a source. If the mechanism is redistribution, particles are not removed; they are moved to lower energy.

A number-conserving cutoff-and-redistribution operator can be represented as
\begin{align}
 f_c(E)&=f_{\mathrm{MJ}}(E)H(E_c-E)+N_{>}g_r(E;E_r,\sigma_r),\\
 N_{>}&=\int_{E_c}^{\infty}f_{\mathrm{MJ}}(E')\,\dd E',\\
 \int_0^\infty g_r(E;E_r,\sigma_r)\,\dd E&=1,
\label{eq:redistribution}
\end{align}
where $g_r$ is a normalized narrow kernel centered at $E_r<E_c$. This model preserves particle number exactly. The change in electron energy density is
\begin{equation}
\Delta U_e=\int_0^\infty E\left[f_c(E)-f_{\mathrm{MJ}}(E)\right]\dd E<0,
\label{eq:deltaU}
\end{equation}
so the missing energy must be transferred to fields, ions, escaping particles, walls, or another reservoir. A self-consistent kinetic simulation must identify that reservoir.

A smoother, more physical stationary distribution results from an energy-dependent loss time,
\begin{equation}
 f(E)\approx \frac{Q(E)+C_{\mathrm{in}}(E)}{\tau_{\mathrm{coll}}^{-1}(E)+\tau_{\mathrm{loss}}^{-1}(E)+\tau_{\mathrm{rad}}^{-1}(E)},
\label{eq:stationary-schematic}
\end{equation}
where $C_{\mathrm{in}}$ represents collisional inflow from adjacent energies. When $\tau_{\mathrm{loss}}$ decreases sharply above a threshold, the high-energy tail is depleted without an artificial discontinuity.

\section{Radiation-weighted moments}

To understand why a small tail can matter, define a generalized radiation-weighted moment
\begin{equation}
M_q[f]=\frac{1}{n_e}\int_0^\infty E^q f_E(E)\,\dd E.
\label{eq:moment}
\end{equation}
If a radiative channel is locally approximated by single-particle power $P(E)\propto E^q$, then the total power is proportional to $n_eM_q$. The exact kernel is not a pure power law, but $M_q$ is a useful diagnostic. A redistribution from $E>E_c$ to $E_r<E_c$ changes the moment by
\begin{equation}
\Delta M_q
=\frac{1}{n_e}\int_{E_c}^{\infty}\left(E_r^q-E^q\right)f_{\mathrm{MJ}}(E)\,\dd E<0
\quad(q>0).
\label{eq:moment-change}
\end{equation}
The inequality is immediate because every moved electron has lower final energy. Its magnitude increases with $q$, explaining why highly energy-weighted radiation is especially sensitive to tail control.

\section{Spectral hardening and biological relevance}

Total-power suppression alone does not reveal which photons were removed. Define the source ratio
\begin{equation}
 R_S(E_\gamma)=\frac{j_{\gamma,c}(E_\gamma)}{j_{\gamma,0}(E_\gamma)}.
\label{eq:spectral-ratio}
\end{equation}
If $R_S$ rises with $E_\gamma$, the residual spectrum is harder even though integrated power falls. Hardening can increase penetration through shielding. If $R_S$ falls with $E_\gamma$, high-energy photons are preferentially removed and distant dose may decrease more strongly than total source power. The shape of $R_S$ is therefore the central input to the transport and biological analysis developed later.

\section{Optical thickness and self-absorption}

The optically thin approximation is valid when
\begin{equation}
\tau_\gamma(E_\gamma)=\int \mu(E_\gamma,s)\,\dd s\ll1
\end{equation}
within the emitting plasma and immediate surroundings. Many magnetically confined plasmas are optically thin to hard X rays, so emitted bremsstrahlung escapes and is a true energy loss. Dense inertial-fusion hot spots can approach conditions where reabsorption and radiation trapping alter the power balance; recent \pB \ analyses emphasize the extreme areal densities required for useful trapping \cite{ochs2025constraints}. The framework therefore allows both regimes, but source suppression has its clearest energy-balance benefit when photons escape.

\section{Diagnostics of the source term}

A validation campaign should combine absolute and relative measurements:
\begin{enumerate}
\item pulse-height or photon-counting detectors for differential spectra;
\item filtered detector arrays or K-edge spectrometers where direct spectroscopy is impractical \cite{chen2008};
\item image plates or radiochromic media for spatially integrated fluence;
\item angularly distributed detectors for anisotropy;
\item electron diagnostics such as Thomson scattering, electron cyclotron emission, hard-X-ray inversion, or magnetic spectrometry; and
\item calorimetry and wall-power accounting to close the energy balance.
\end{enumerate}
Bremsstrahlung inversion is ill posed, so detector response, pileup, regularization, and uncertainty must be included \cite{voss1992,kontar2011,swanson2018,meadowcroft2012}.

\section{Summary}

Relativistic plasma radiation is controlled by the entire electron distribution, not only by a fitted temperature. Electron--ion and electron--electron bremsstrahlung, synchrotron emission, and secondary photon processes respond differently to energy, pitch angle, density, composition, and magnetic field. A cutoff can reduce radiation only if its particle and energy balances are physically closed. The next chapter develops magnetic mechanisms capable of producing the required energy-selective transport.
